\chapter{Materials and Methods}
\label{ch:methods}

This is a theoretical and computational thesis: its "materials" are the
published SME and Dobrescu--Mocioiu (DM) formalisms reviewed in Chapter~2
and the compiled experimental constraint literature, and its principal
research tool is SPINDEP, a purpose-built analysis pipeline (Python,
NumPy/SciPy, and symbolic computation in SymPy) developed for this thesis.
This chapter describes the research procedures used to derive the
SME$\to$DM coefficient mapping (Section~\ref{sec:procedure}) and the
methods of analysis used to compile and statistically compare experimental
constraints (Section~\ref{sec:analysis-methods}). The results these
procedures produce are presented and discussed in Chapter~4.

\section{Research Procedures: Foldy--Wouthuysen Reduction of SME Coefficients}
\label{sec:procedure}

Each SME fermion-sector coefficient of interest, $b_\mu$, $H_{\mu\nu}$, and
$d_{\mu\nu}$, was reduced to its non-relativistic Hamiltonian and matched
onto the Dobrescu--Mocioiu potential catalogue by the same five-step
procedure.

\begin{enumerate}
  \item \textbf{Write the covariant Lagrangian term.} The relevant term is
    isolated from the general minimal SME Lagrangian
    (Eq.~\ref{eq:sme-minimal}), and its CPT and Lorentz transformation
    properties are checked directly from the transformation of the
    corresponding fermion bilinear (or, for the derivative-coupled
    coefficient $d_{\mu\nu}$, stated following Kostelecký and Lane, 1999,
    since the bilinear rule alone does not apply to a derivative-coupled
    operator).
  \item \textbf{Derive the correction to the Dirac Hamiltonian.} Varying the
    modified Lagrangian and solving the resulting modified Dirac equation
    for $H\psi=E\psi$ gives the coefficient's contribution to the free Dirac
    Hamiltonian $H_0=\boldsymbol\alpha\cdot\mathbf p+\beta m$, in the Dirac
    representation of the gamma matrices ($\gamma^0=\mathrm{diag}(1,1,-1,-1)$,
    $\gamma_5$ off-diagonal).
  \item \textbf{Classify each contribution as even or odd.} Using standard
    Dirac-algebra identities ($\{\gamma_5,\gamma^\mu\}=0$ and the explicit
    $4\times4$ matrix forms of $\gamma^0\gamma_5\gamma^i$,
    $\gamma^0\sigma^{ij}$, and related combinations, verified symbolically
    with a purpose-written Dirac-algebra module), each component of the
    correction is classified as block-diagonal (\emph{even}, commuting with
    $\beta$) or block-off-diagonal (\emph{odd}, anticommuting with $\beta$).
  \item \textbf{Apply the Foldy--Wouthuysen expansion.} An even component
    contributes to the non-relativistic Hamiltonian unchanged at order
    $m^0$, requiring no further transformation. An odd component is
    combined with the leading odd operator $\boldsymbol\alpha\cdot\mathbf p$
    through the general order-$1/m$ FW expansion (Eq.~\ref{eq:FWexpansion}),
    producing a subleading, momentum-dependent contribution. This procedure
    is preferred over a direct equation-of-motion expansion in the full
    four-component spinor because it delivers the particle-sector
    Hamiltonian as an explicit $2\times2$ Pauli expression whose operator
    structure can be matched term-by-term against the closed-form $V_i$
    potentials of Eqs.~\ref{eq:v1}--\ref{eq:v16}. For $b_\mu$,
    this step is additionally cross-checked against the antiparticle
    sector directly from the block structure identified in step 3: since
    the even/odd classification already block-diagonalises the operator,
    the negative-energy (lower) block gives the antiparticle's matrix
    element in the original, unconjugated Hamiltonian, and the standard
    hole-theory sign flip ($H\to-H$ on reinterpreting a negative-energy
    solution as a positive-energy antiparticle) converts it into the
    antiparticle's non-relativistic Hamiltonian. This route was adopted in
    preference to a direct charge-conjugation argument on the bilinear
    $\bar\psi\gamma_5\gamma^\mu\psi$ after that bilinear was verified, with
    the same Dirac-algebra tooling, to be even under $C$ for every $\mu$
    and therefore unable by itself to produce the required sign
    difference; the resulting sign flip (or its absence) is nonetheless
    the mechanism responsible for a genuine matter--antimatter asymmetry.
  \item \textbf{Match to the Dobrescu--Mocioiu catalogue.} The resulting
    single-particle non-relativistic Hamiltonian is promoted to a two-body
    potential by pairing the same vertex structure on both interacting
    fermions, and the resulting operator structure (in $\boldsymbol\sigma_1$,
    $\boldsymbol\sigma_2$, $\hat{\mathbf r}$, and $\mathbf v$) is compared
    directly against the closed forms of $V_1$--$V_{16}$
    (Eqs.~\ref{eq:v1}--\ref{eq:v16}) to identify which potential(s) the
    coefficient generates and at what order in $1/m$.
\end{enumerate}

Every derivation obtained by this procedure was additionally verified by an
independent, executable symbolic computation (SymPy notebooks
\texttt{FW\_bmu}\allowbreak\texttt{\_term.ipynb},
\texttt{FW\_Hmunu}\allowbreak\texttt{\_term.ipynb}, and
\texttt{FW\_dmunu}\allowbreak\texttt{\_term.ipynb}), including, for the CPT asymmetry parameter
specifically, a symbolic limit calculation used to check the behaviour of
Eq.~\eqref{eq:asymmetry} under an exact signed CPT-odd substitution
(Chapter~4, Section~4.1.3). Where a derived result could be checked against
an independently published non-relativistic Hamiltonian (Kostelecký and
Lane, 1999, Eq.~4), this comparison was carried out component by component
rather than only at the level of overall structure, which is how the open
sign in the $d_{ij}$ and $d_{00}$ components (Chapter~4) was identified
rather than overlooked.

\section{Methods of Analysis}
\label{sec:analysis-methods}

\subsection{Constraint Database Compilation}
\label{sec:db-methods}

Each experimental constraint used in this thesis is stored as a
two-column comma-separated value (CSV) file (interaction range $\lambda$ in
metres, upper bound on the coupling $|g|$), organised in a directory tree
by coupling family and interaction class, following the naming convention
\texttt{\{V\}\{Author\}\{Year\}}\allowbreak\texttt{\_m\_abs\_}\allowbreak\texttt{\{sector\}.csv}. A parser
recursively discovers every CSV file under the dataset root and extracts
the potential label, source citation, and fermion sector from the filename
and directory path, resolving alternate spellings of antimatter sectors
(for example \texttt{ebare}, \texttt{ebar}, and \texttt{epbare} all resolve
to the canonical labels \texttt{eebar} and \texttt{epbar}). Files whose
sector or potential token cannot be resolved are retained in the compiled
registry with the corresponding field recorded as \texttt{UNKNOWN} rather
than being silently dropped, so that any classification gap remains visible
rather than hidden by omission. Such files are then held out of the physics
analysis by an explicit, per-file exclusion list that records a stated
reason for each, rather than by a blanket rule on the \texttt{UNKNOWN}
label; the registry therefore reports both what was compiled and what was
analysed, and the difference is auditable. Interaction ranges
reported in the literature in a mix of units (metres, MeV$^{-1}$,
eV$^{-1}$, centimetres, nanometres, femtometres) are converted to a common
metre scale by a unit-detection routine that inspects each filename for an
explicit unit token before the dataset is used in any downstream analysis.
Two datasets are matched into a matter--antimatter pair only when they
share the same coupling family, the same DM potential label, and the same
interaction class, and their fermion sectors are physically conjugate (for
example \texttt{ee} with \texttt{eebar}), with exactly one of the two
flagged as an antimatter dataset; this pairing logic is deliberately
conservative and will not match two datasets differing in potential
classification even where their sectors are conjugate.

\subsection{Statistical Methodology for the Asymmetry Analysis}
\label{sec:stat-methods}

For each matched matter--antimatter pair, both constraint curves are
interpolated (log-linear in $\log\lambda$--$\log g$) onto a common
log-spaced grid of 300 points spanning their overlapping $\lambda$ range,
and the asymmetry parameter $A_\alpha(\lambda)$ of Eq.~\eqref{eq:asymmetry}
is evaluated at every grid point. Statistical significance is assessed in
four stages, each addressing a limitation of the one before it.

\begin{enumerate}
  \item A \emph{uniform $\chi^2$ test} assumes a flat 10\% fractional
    uncertainty at every grid point.
  \item A \emph{weighted $\chi^2$ test} instead estimates a per-point
    fractional uncertainty from the local curvature of each constraint
    curve in log-log space, on the grounds that a smoothly varying
    published bound is better determined than one with sharp local
    features.
  \item Both tests above treat the 300 grid points as independent degrees
    of freedom, which is not justified, since the points are interpolated
    from a much smaller number of real measurements along smooth curves and
    adjacent points are strongly correlated. An \emph{effective degrees of
    freedom} is therefore estimated from the autocorrelation length of the
    $\chi^2$ residuals in grid-index space (the lag at which the residual
    autocorrelation first drops below $1/e$), and the weighted-$\chi^2$
    significance is recomputed against this corrected, typically much
    smaller, degrees of freedom.
  \item A \emph{parametric bootstrap} (2000 resamples, propagating the
    estimated per-point uncertainty through to $A_\alpha$ itself) gives a
    95\% confidence interval on the mean $|A_\alpha|$ for each pair.
\end{enumerate}

Two properties of this bootstrap step are worth stating explicitly rather
than leaving implicit in the word ``parametric'' above. First, it is a
delta-method approximation, not a resampling of the underlying raw
matter/antimatter coupling values: by the point the bootstrap runs, each
pair has already been reduced to $A_\alpha(\lambda)$ and its estimated
combined uncertainty, so the per-point perturbation used to generate each
resample is obtained by linearising $A_\alpha=(g_m-g_a)/(g_m+g_a)$ about
$g_m\approx g_a$ and propagating the known uncertainty through that
linearisation, rather than by redrawing $(g_m,g_a)$ pairs with
replacement. Second, because the confidence interval is computed on
$|A_\alpha|$ rather than the signed asymmetry, it is systematically
biased away from zero even for a pair with no true asymmetry: folding a
distribution centred at zero through an absolute value can only push
resampled values upward, never down, so a bootstrap confidence interval
that excludes zero should not, on its own, be read as evidence of a
non-zero asymmetry.

This methodology is applied uniformly to every matched pair in the
compiled database (Chapter~4), and its results are interpreted throughout
in light of the analytic sensitivity-gap argument developed in
Section~\ref{sec:asymmetry-def} and Chapter~4: statistical significance of
an asymmetry is a separate claim from evidence of CPT violation, and the
methodology above establishes the former without, on its own, establishing
the latter.
